% source: https://github.com/pfradin/luadraw/discussions/301#discussioncomment-17976473 (pfradin, block 1)
\documentclass[12pt]{standalone}
\usepackage[svgnames]{xcolor}
\usepackage[3d]{luadraw}
\begin{document}
\begin{luadraw}{name=tetrahedron1} 
local ld = luadraw 
local pt3d = ld.pt3d 
local M, Mc, O, vecK = pt3d.M, pt3d.Mc, pt3d.Origin, pt3d.vecK 
local sin, cos, pi, sqrt = math.sin, math.cos, math.pi, math.sqrt 
local g = ld.graph3d:new{window={-2,6,-1,5}, viewdir={240,60} ,size={12,12}} 
ld.Hiddenlinestyle = "dashed"; 
local ASB,BSC,ASC=60,60,45 
local as,bs,cs = 5, 6, 7
local ab,bc,ac= sqrt(as^2 + bs^2 - 2*as*bs*cos(ASB*ld.deg)),sqrt(bs^2+cs^2-2*bs*cs*cos(BSC*ld.deg)), sqrt(as^2+cs^2-2*as*cs*cos(ASC*ld.deg)) 
local A,B,C,S = ld.tetra_len(ab,ac,as,bc,bs,cs) 
local T = ld.tetra(A,B-A,C-A,S-A) 
local H = ld.proj3d(A,{S,pt3d.prod(B-S,C-S)})
local u = pt3d.prod(C-S,A-S)
local v = pt3d.prod(B-S,C-S)
local a = pt3d.angle3d(u,v)
if a > math.pi/2 then a = math.pi-a end -- to obtain the acute angle <- change here
local I = (3*C+S)/4 -- a dot on (SC)
local d1, d2 =  pt3d.prod(u,C-S), pt3d.prod(C-S,v)
-- a point J1 in (SCA) with (IJ1) and (IC) orthogonal
-- a point J2 in (SCB) with (IJ2) and (IC) orthogonal
local J1, J2 = I + 0.7*pt3d.normalize(d1), I + 0.7*pt3d.normalize(d2)
g:Dpoly(T, {color="Crimson",mode=ld.mShadedHidden}) 
g:Ddots3d({A,B,C,S,H}) 
g:Dpolyline3d({A,H},"dashed")
g:Dangle3d(A,H,B,0.25,"dashed") 
-- angle a
g:Dpolyline3d({I,J1}); g:Dpolyline3d({I,J2},"dashed")
g:Dangle3d(C,I,J1,0.25); g:Dangle3d(C,I,J2,0.25,"dashed")
g:Darc3d(J1,I,J2,0.5,1,"dashed")
g:Labelsize("footnotesize")
g:Dlabel3d( "$A$",A,{pos="S"}, "$B$",B,{pos="S"}, "$C$",C,{pos="N"},"$S$",S,{pos="N"},
    "$"..ld.strReal(a*ld.rad,2).."^\\circ$",(J1+J2)/2,{pos="E"},"$H$",H,{pos="N"}) 
g:Show() 
\end{luadraw}
\end{document} 
